\documentclass[landscape]{article}
\usepackage[pdftex]{graphicx}
\usepackage[utf8]{inputenc}
\usepackage{enumerate}
\usepackage{icomma}
\usepackage{amssymb}
\usepackage{tikz}
\usepackage{href-ul}
\hypersetup{
	colorlinks=true,
	linkcolor=blue,
	urlcolor=blue}
\usepackage{geometry}
\geometry{a4paper, top=15mm, left=15mm, right=15mm, bottom=15mm,
	headsep=10mm, footskip=12mm}
\twocolumn

\begin{document}
	\begin{tikzpicture}(10,3)
		\draw[ultra thick](2,0) --(10,0) -- (10,1.3) --(2,1.3) -- (2,0);
		\draw[fill=black](2,0)-- (0,.6) -- (2,.6) -- (2,0);
		\node at (6,.6) {\large Solution sheet from \href{http://www.okuyakl.com}{www.okuyakl.com}};
	\end{tikzpicture}
	\vspace{0.5 cm}
	
	\noindent{\bf Task 1. a)}\\
	$$
	\renewcommand{\arraystretch}{2}
	\begin{array}{rcll}
		\sin(\alpha + 60^\circ) &=& {3 \over 6} & | \sin^{-1} \\
		\alpha_1 + 60^\circ &=& 30^\circ \quad (*) & | - 60^\circ \\
		\alpha_1 &=& -30^\circ & | +360^\circ \\
		\alpha_1 &=& 330^\circ \\
		(*)\quad 180^\circ -(\alpha_2+ 60^\circ) &=& 30^\circ \\
		120^{\circ}-\alpha_2 &=& 30^\circ \\
		\alpha_2 &=& 90^\circ
	\end{array}\\
	\mathbb{L}=\{90^\circ ;330^\circ \}
	$$
	\noindent{\bf Task 1. b)}\\
	$$
	\renewcommand{\arraystretch}{2}
	\begin{array}{rcll}
		\cos(\alpha +45^\circ)&=&{2 \over 2} & |\cos^{-1}\\
		\alpha_1 +45^\circ &=& 0^\circ &\quad (*)\\
		\alpha_1 &=& -45^\circ &| +360^\circ \\
		\alpha_1 &=& 315^\circ \\
		(*) \quad 360^\circ - (\alpha_2+45^\circ) &=& 0^\circ \\
		315^\circ - \alpha_2 &=& 0^\circ \\
		\alpha_2 &=& 315^\circ &=\alpha_1 \\
	\end{array}\\
	\mathbb{L}=\{ 315^\circ \}
	$$
	
	\newpage
	\noindent{\bf Task 1. c)}\\
	$$
	\renewcommand{\arraystretch}{2}
	\begin{array}{rcll}
		\sin^2\alpha &=&0.75& |\pm \sqrt{\quad} \\
		\sin\alpha &=& \sqrt{0.75}&|\sin^{-1} \\
		\alpha_1 &=& 60^\circ \\
		\alpha_2 &=& 180^\circ - 60^\circ&= 120^\circ\\
		\alpha_3 &=& -60^\circ &| +360^\circ\\
		\alpha_3 &=& 300^\circ\\
		\alpha_4 &=& 180^\circ-(-60^\circ)&=240^\circ\\
	\end{array}\\
	\mathbb{L}=\{60^\circ;120^\circ; 240^\circ; 300^\circ \}
	$$
	
	\noindent{\bf Task 1. d)}\\
	$$ \sin\alpha + \cos\alpha = 2$$
	There is no angle for which this applies because sine and cosine are 1 in different places.
	$$\mathbb{L}=\{ \}$$
	
	\noindent{\bf Task 1. e)}\\
	$$
	\renewcommand{\arraystretch}{2}
	\begin{array}{rcll}
		1-\cos^2\alpha &=& 0\\
		\cos^2\alpha &=& 1 & \pm \sqrt{\quad}\\
		\cos\alpha_1 &=& 1 \\
		\alpha_1 &=& 0^\circ\\
		\alpha_2 &=& 360^\circ\\
		\cos\alpha_3 &=& -1 \\
		\alpha_3 &=& 180^\circ\\
	\end{array}\\
	\mathbb{L}=\{0^\circ; 180^\circ ; 360^\circ \}
	$$
	
	\newpage
	\noindent{\bf Task 1. f)}\\
	We apply the third binomial formula:
	$$
	\renewcommand{\arraystretch}{2}
	\begin{array}{rcll}
		\sin^2\alpha -0.36 &=& 0.28 \\
		\sin^2\alpha &=& 0.64 &|\sqrt{\quad}\\
		\sin\alpha &=& \pm 0.8 \\
		\alpha_1 &=& 53.13^\circ \\
		\alpha_2 &=& -53.13^\circ &= 306.87^\circ \\
		\alpha_3 &=& 180^\circ-53.13^\circ &= 126.87^\circ\\
		\alpha_4 &=& 180^\circ-(-53.13^\circ) &= 233.13^\circ\\
	\end{array}\\
	\mathbb{L}=\{53,13^\circ; 126.87^\circ; 233.13^\circ \}
	$$
	
	\noindent{\bf Task 1. g)}\\
	From this task onwards we have to use the addition theorems; from one follows:
	$\sin(2\alpha)=2\cdot\sin\alpha\cdot\cos\alpha$
	$$
	\renewcommand{\arraystretch}{2}
	\begin{array}{rcll}
		\sin(2\alpha)&=&0.5\\
		2\alpha_1 &=& 30^\circ \\
		\alpha_1 &=& 15^\circ \\
		180^\circ-2\alpha_2 &=& 30^\circ\\
		2\alpha_2 &=& 150^\circ \\
		\alpha_2 &=& 75^\circ
	\end{array}\\
	\mathbb{L}=\{15^\circ; 75^\circ \}
	$$
	
	\newpage
	\noindent{\bf Task 1. h)}\\
	From the addition theorems it also follows: $\cos(2\alpha)=\cos^2\alpha-\sin^2\alpha$
	$$
	\renewcommand{\arraystretch}{2}
	\begin{array}{rcll}
		\cos\alpha &=& \sin\alpha& |:\cos{\alpha}\\
		1&=&{ \sin\alpha \over \cos\alpha}\\
		1&=&\tan\alpha& |\tan^{-1}\\
		\alpha_1 &=& 45^\circ\\
		\alpha_2 &=&180^\circ+45^\circ&= 225^\circ
	\end{array}\\
	\mathbb{L}=\{45^\circ; 225^\circ \}
	$$
	
	\noindent{\bf Task 1. i)}\\
	We have a quadratic equation and substitute with $u=\sin\alpha$:
	$$4u^2 + 4u +1 =0$$
	The midnight formula gives:
	$$u_{1/2}=0.5$$
	We can back-substitute:
	$$\sin\alpha={0.5} \Rightarrow \alpha_1 =30^\circ; \alpha_2=150^\circ$$
	$$\mathbb{L}=\{30^\circ;150^\circ \}$$
	
	\noindent{\bf task 1. j)}\\
	We divide by 2:
	$$\renewcommand{\arraystretch}{2}
	\begin{array}{rcll}
		\sin{\alpha\over 2} &=& 0\\
		{\alpha_1\over 2} &=& 0 \\
		\alpha_1 &=& 0\\
		{\alpha_2 \over 2} &=&180^\circ \\
		\alpha_2 &=& 360^\circ &= 0^\circ\\
		{\alpha_3 \over 2} &=&360^\circ \\
		\alpha_3 &=& 720^\circ &=0^\circ \\
	\end{array}\\
	\mathbb{L}=\{0^\circ; 360^\circ \}
	$$
	
	\newpage
	\noindent{\bf Task 1. k)}\\
	Same numerators on the left and right, so we can directly equate the denominators:
	$$
	\renewcommand{\arraystretch}{2}
	\begin{array}{rcll}
		4 \sin\gamma&=& 3\\
		\sin(\gamma) &=& {3\over 4}\\
		\gamma_1 &=& 48.59^\circ \\
		\gamma_2) &=& 180^\circ - 48.59^\circ \\
		\gamma_2 &=& 131.41^\circ
	\end{array}\\
	\mathbb{L}=\{48.59^\circ; 131.41^\circ \}
	$$
	
\noindent{\bf Task 1. l)}\\
We switch and apply an addition theorem:
$$
\renewcommand{\arraystretch}{2}
\begin{array}{rcll}
	\sin\epsilon &=& 4 \sqrt{3}\cdot (\sin\epsilon \cos(30^\circ)+\sin(30^\circ)\cos\epsilon)\\
	\sin\epsilon &=& 4 \sqrt{3}\cdot (\sin\epsilon {1\over 2}\sqrt{3} + {1\over 2}\cos\epsilon)\\
	\sin\epsilon &=& 2\cdot 3 \cdot \sin\epsilon +2 \sqrt{3} \cos\epsilon\\
	-5 \sin\epsilon &=& 2 \sqrt{3} \cos\epsilon &|:\cos\epsilon \quad|:(-5)\\
	\tan\epsilon &=& -{2\over 5}\sqrt{3}\\
	\epsilon_1 &=& -34.71^\circ\\
	\epsilon_1 &=& 325.29^\circ\\
	\epsilon_2 &=& 180^\circ-34.71^\circ = 145.29^\circ\\
\end{array}\\
\mathbb{L}=\{145.29^\circ; 325.29^\circ \}
$$

\begin{center}
	\includegraphics[width=7 cm]{../../viecher/eendcomic.pdf}
	
	Here you can go back to the \href{https://www.okuyakl.de/math/m10gonL082/ae082.pdf}{task sheet}
\end{center}


\end{document}